\section{Related Work}
\label{sec:related}

\para{Model collapse.}
\citet{shumailov2023curse} introduced the term \emph{model collapse} and demonstrated it across VAEs, Gaussian mixture models, and language models (OPT-125M on \wikitext). They identified two error sources---statistical approximation error from finite sampling and functional approximation error from limited model expressiveness---and noted that model-generated data produces ``very small gradients.'' Our work makes this gradient observation precise and shows that the directional structure of gradients, not just their magnitude, is the critical factor. \citet{wang2024theoretical} provided the first theoretical proof that autoregressive language models must collapse when trained recursively on finite generated data, establishing inevitability regardless of the rate at which synthetic data enters the corpus. \citet{dohmatob2024strong} showed that even 1\% synthetic data contamination can trigger collapse in simplified models, underscoring the sensitivity of the phenomenon.

\para{Distributional analysis of collapse.}
\citet{dohmatob2024tale} connected model collapse to neural scaling laws through tail truncation, showing that synthetic data produced by top-$p$ sampling and temperature scaling systematically removes low-probability tokens. They derived double and triplet scaling laws for synthetic data regimes. \citet{kazdan2024bad} provided a systematic statistical analysis under Replace and Accumulate paradigms, finding that data accumulation mitigates but does not eliminate collapse. \citet{feng2024synthesize} proposed verification-based curation to maintain synthetic data quality. \citet{seddik2025position} challenged the prevailing narrative, showing that mixing real data effectively prevents the worst outcomes. Our gradient-direction analysis provides the optimization-level mechanism underlying these distributional observations: tail truncation in data space manifests as gradient signal loss in tail-representing parameter directions.

\para{Gradient and weight update analysis.}
The closest work to ours is \citet{wu2025mitigating}, who measured weight update norms ($\|\Delta \mW\|$) during fine-tuning and found that self-generated data produces roughly $3\times$ smaller updates than ground-truth data. They introduced Selective Token Masking to mitigate forgetting by masking high-perplexity tokens. However, their analysis focused on aggregate weight norms rather than per-step gradient dynamics or directional decomposition. \citet{wang2025synthetic} proposed generating synthetic text that matches gradient signals from real data, directly addressing the gradient mismatch we characterize. Our work complements these efforts by providing the first fine-grained analysis of gradient directions, revealing that the problem is not merely one of magnitude but of dimensionality.

\para{Synthetic data quality and curation.}
\citet{el2024maximizing} used random matrix theory to model classifier performance on mixed real/synthetic data, finding smooth phase transitions as synthetic data quality varies. Their framework accounts for both feature distribution shifts and label noise, providing a mathematical basis for understanding real/synthetic data mixtures in high dimensions. More broadly, the growing literature on synthetic data curation~\citep{feng2024synthesize,seddik2025position} focuses on data-level interventions. Our gradient-level analysis suggests that effective curation strategies should aim to maximize gradient subspace coverage, not just data diversity as measured by surface statistics.
